% Fichier complété par tools/complete_missing_corrections.py.
% Source : CIN/CIN-03-Transmetteurs/93_Lokomat/93_Lokomat.tex
% Statut : rédigé (2 question(s)).
\begin{corrigebox}[Corrigé rédigé]
\CorrigeQuestion{1}
On applique la relation de Willis au train épicycloïdal du premier
                étage en tenant compte de l'élément fixe indiqué sur le schéma. Avec les
                dentures données, le rapport de réduction de l'étage vaut
                \[r_1=\frac{\omega_{\rm sortie}}{\omega_{\rm entrée}}
                =\frac{Z_1Z_3}{Z_0Z_2}=\frac{18\times24}{60\times45}
                =\boxed{0,16}\]
                (le signe dépend du nombre d'engrènements extérieurs).
\CorrigeQuestion{2}
Les trois étages étant identiques et montés en cascade,
                $r_g=r_1^3=0,16^3=\boxed{4,096\times10^{-3}}$, soit un rapport de
                réduction en vitesse voisin de $1/244$.
\end{corrigebox}
